\section{Equality read from one fixed row}
\label{sec:first-row-mathematics}

The new reduction is local to witness discharge.  It removes the equality
proof that the preceding \emph{Name} rule asks the certificate to supply.
Instead, the rule reads one specified instance of its own opened clause.
The zero instance must be the fixed reflexive equation $0=0$, and the
successor instance must name the fresh witness.  These are decidable
syntactic tests, not proof-search operations.

\subsection{Setting and complete discharge rule}

Write $N$ for full computation normalization of scalar iterator terms and
of the numerical terms in formulas.  The reductions are
\[
 I(0,a;u.h)\longrightarrow a,
 \qquad
 I(Sn,a;u.h)\longrightarrow h[I(n,a;u.h)/u].
\]
Terms and formulas in inference labels are normal forms.  Substitution is
capture avoiding and followed by $N$.  In particular, the argument uses
the inherited normalization--substitution identity
\begin{equation}
 N\bigl(N(F)[t/b]\bigr)=N\bigl(F[t/b]\bigr).
 \label{eq:row-normalization-coherence}
\end{equation}
This is a metasyntactic identity.  It is not an application of an
object-language equality rule.  No numerical variable ranges over
programs or function descriptions.

Before normalization the complete numerical grammar is
\[t::=x\mid 0\mid S(t)\mid I(n,a;u.h),\]
where only $u$ is bound in $h$, and the other arguments do not bind.
The native formulas are equations and $D_{x,y}B$.  Its ordinary reading is
\[
 \bigl(D_{x,y}B\bigr)^\circ
   =\exists x\bigl(B^\circ(x,0)\to\forall y\,B^\circ(x,y)\bigr).
\]
A context consists of reusable clauses $[\Pi;y.A(y)]$, with zero or one
antecedent in $\Pi$; the displayed $y$ is not free in the antecedent.
An ordinary assumption $A$ is the unconditional constant clause
$[\,;v.A]$, where $v$ is unused.  Every derivation has the fixed initial
clause
\[
                    \mathcal R=[\,;y.y=y].
\]
Its instances supply reflexivity by the ordinary clause-use instruction.
The designated contradiction is still $0=S0$.

The Use instruction is unchanged: from the antecedent of a selected
clause, if present, infer its consequent at any supplied numerical term.
The Intro instruction is also unchanged.  At witness $t$ and fresh $z$,
it derives $D_{x,y}B$ from a branch proving $N(B(t,Sz))$.  The branch has
the ordinary assumption $N(B(t,z))$ precisely when $B$ is an equation and
$t$ is literally $0$; otherwise it has $N(B(t,0))$.  The witness and the
matrix's passive parameters are outer scoped, and $z$ is fresh for all
retained data.

For the new discharge instruction, first prove $D_{x,y}B$ and open it at
a fresh numerical parameter $b$.  Write
\[
 \mathcal C_b=[N(B(b,0));y.N(B(b,y))],\qquad 1=S0.
\]
Define the term $\tau_B(b)$ deterministically:
\begin{equation}
 \tau_B(b)=
 \begin{cases}
 t,&N(B(b,0))\text{ is literally }0=0
       \text{ and }N(B(b,1))\text{ is literally }b=t,\\
 b,&\text{otherwise}.
 \end{cases}
 \label{eq:row-target}
\end{equation}
There is no certificate for either test.  The fresh witness must be the
\emph{entire left term} of the tested equation, not merely one of its
subterms.  The rule is
\begin{equation}
 \frac{\Gamma\vdash D_{x,y}B\qquad
       \Gamma,\mathcal C_b\vdash F}
      {\Gamma\vdash N(F[\tau_B(b)/b])},
 \qquad b\notin\operatorname{FV}\bigl(N(F[\tau_B(b)/b])\bigr).
 \tag{Row}
 \label{rule:first-row}
\end{equation}
The parameter $b$ is fresh for the retained context and the matrix's
passive parameters.  Replacement is simultaneous and performed once:
inserted copies of $\tau_B(b)$ are not revisited.  The selected target may
itself contain $b$.  The condition required for discharge is absence of
$b$ from the \emph{final normalized result}, not absence from that target.

Use, Intro, and Row, together with $\mathcal R$ and the specified canonical
syntax, are the entire presentation.  Full normalization is part of its
trusted specification and implementation.  There are no additional
primitive equality, computation, successor, or induction rules.  This
still counts as three instruction families; the reduction is from three
proof premises to two in the discharge instruction and the deletion of
its independently supplied equality data.

\subsection{Soundness and the equality-free preliminaries}

Let $\mathcal K_{=}$ denote the preceding full-normal Name calculus,
whose discharge is the explicit rule
\[
 \frac{\Gamma\vdash D_{x,y}B\qquad\Gamma,\mathcal C_b\vdash b=t
       \qquad\Gamma,\mathcal C_b\vdash F}
      {\Gamma\vdash N(F[t/b])},
 \quad b\notin\operatorname{FV}(N(F[t/b])).
 \tag{Name}
\]
Here $b$ has the same freshness requirement as in Row, and $t$ is
well scoped in the local context; it may contain $b$.  Let $\mathcal K_{\mathrm{row}}$ denote the new calculus.

\begin{lemma}[The rule supplies its own equality]
\label{lem:row-selected-equality}
In the opened context $\Gamma,\mathcal C_b$, one can derive
$b=\tau_B(b)$ using only Use and $\mathcal R$.
\end{lemma}
\begin{proof}
If the test in \eqref{eq:row-target} succeeds, reflexivity gives its zero
antecedent $0=0$, and Use of $\mathcal C_b$ at $1$ gives exactly $b=t$.
If the test fails, Use of $\mathcal R$ at $b$ gives $b=b$.
\end{proof}

\begin{corollary}[Embedding into the preceding calculus]
\label{cor:row-soundness}
Every Row inference is derivable in $\mathcal K_{=}$ with the same end
sequent.  Hence every new derivation is sound for the preceding ordinary
interpretation.
\end{corollary}
\begin{proof}
Retain the major and formula branches, insert the equality proof of
Lemma~\ref{lem:row-selected-equality}, and use the old Name inference.
Its result and eigenparameter condition are exactly those of Row.
\end{proof}

The same argument works directly in ordinary intuitionistic equality
logic: equality gives $F[\tau_B(b)/b]$, normalization preserves its
meaning, and the result's $b$-freedom permits existential elimination.
This also explains why no occurs check on the selected target is needed.

\begin{lemma}[Ordinary paired elimination]
\label{lem:row-ordinary-elimination}
From $\Gamma\vdash D_{x,y}B$ and
$\Gamma,\mathcal C_b\vdash A$, with $b$ absent from $A$, one can derive
$\Gamma\vdash A$.
\end{lemma}
\begin{proof}
Use Row with the given formula branch $A$.  Regardless of which target
the test selects, $N(A[\tau_B(b)/b])$ is literally $A$: there is no $b$ to
replace, and $A$ is already normal.  In particular, this argument does
not assume that the rule takes its fallback branch.
\end{proof}

Weakening, injective fresh renaming, and cut of an ordinary assumption
are admissible by direct induction on derivations.  Freshen local
parameters before insertion.  An ordinary-assumption cut replaces uses
of its unconditional constant clause by the supplied proof.  The matrix
and deterministic test at a Row node are unchanged by such insertion.
Injective renaming commutes with normalization and preserves and reflects
the tests against $0=0$ and against the distinguished fresh parameter.
These structural arguments do not assume numerical substitution of
derivations or any equality-congruence theorem.

For a variable $w$ unused in $A$, put
$q_yA(y)=D_{w,y}A(y)$.

\begin{lemma}[Numerical cases before equality transport]
\label{lem:row-cases}
The following rules are derivable using only Use, Intro,
Lemma~\ref{lem:row-ordinary-elimination}, and $\mathcal R$:
\begin{align*}
 &\frac{\Gamma,A(0)\vdash A(Sz)}
        {\Gamma\vdash q_yA(y)},\\[4pt]
 &\frac{\Gamma\vdash q_yA(y)\qquad\Gamma\vdash A(0)}
        {\Gamma\vdash A(t)},\\[4pt]
 &\frac{\Gamma\vdash A(0)\qquad\Gamma\vdash A(Sz)}
        {\Gamma\vdash A(t)}.
\end{align*}
Here all instances are normalized, and the successor-branch parameter is
fresh for the retained context and the template's passive parameters.
\end{lemma}
\begin{proof}
For the first rule, introduce the unused witness $w$ at $1$.  This forces
Intro's case mode even when $A$ is an equation, because the chosen
witness is not literally $0$.  Its branch is precisely the supplied
$A(0)\vdash A(Sz)$.

For the second rule, open $q_yA$ by ordinary paired elimination.  Since
its witness coordinate is unused, the opened clause is
$[A(0);y.A(y)]$.  Use it at $t$ with the given proof of $A(0)$, then
discharge the unused witness.  For the third rule, weaken its successor
proof by the unused assumption $A(0)$ and combine the first two rules.
\end{proof}

This is numerical case distinction, not equation induction: its
successor branch has no induction hypothesis.  It has been obtained
before, and without, any transport principle.

\subsection{Recovering arbitrary equality from two numerical rows}

Define a normalizing selector and a short parity term by
\[
 C(v,a,c)=N(I(v,a;u.c)),\qquad
 P(y)=N(I(y,0;u.C(u,1,0))),
\]
where the state variable in the first expression is unused in $c$.
Only four literal computation identities are needed:
\begin{equation}
 C(0,a,c)=a,\quad C(Sv,a,c)=c,\quad
 P(1)=1,\quad P(2)=0,\qquad 2=S1.
 \label{eq:row-endpoints}
\end{equation}
No predecessor, parity range theorem, or arithmetic induction occurs in
these finite computations.

The following construction adapts the two-row witness package in the
preceding witness-discharge report~\cite{prior}.  Its rows are
exchanged: row $1$ now contains the equation the rule must read
automatically, and row $2$ contains the formula the branch supplies.
The use of the selector and parity term is inherited; the fixed-row
discharge restriction is the present change.

\begin{theorem}[Uncrossed term transport]
\label{thm:row-term-transport}
For every fixed numerical term template $h(v)$, the new calculus derives
\[
 \frac{\Gamma\vdash s=t\qquad\Gamma\vdash N(h(s))=r}
      {\Gamma\vdash N(h(t))=r}.
\]
The retained context may contain arbitrary clauses.
\end{theorem}
\begin{proof}
Choose fresh matrix variables $x,y$ and define the equation
\[
 B(x,y)=\bigl(
 C(y,0,C(P(y),N(h(x)),x))
 =C(y,0,C(P(y),r,t))\bigr).
\]
All displayed expressions are normalized.  By
\eqref{eq:row-endpoints}, its relevant instances are literally
\begin{equation}
 B(x,0)=(0=0),\qquad B(x,1)=(x=t),\qquad
 B(x,2)=(N(h(x))=r).
 \label{eq:row-three-instances}
\end{equation}

We first prove $D_{x,y}B$ by Intro at the witness $s$.  For its successor
branch consider the formula
\[
 A(v)=\bigl(C(v,N(h(s)),s)=C(v,r,t)\bigr).
\]
Its zero instance is the supplied equation $N(h(s))=r$.  Every successor
instance is the supplied equation $s=t$.  Applying
Lemma~\ref{lem:row-cases} at the target $P(Sz)$ gives
\[
 C(P(Sz),N(h(s)),s)=C(P(Sz),r,t),
\]
which is exactly $N(B(s,Sz))$.  The eigenparameter for this case proof is
chosen fresh independently of $z$.

This successor proof does not use Intro's local hypothesis.  It is
therefore valid both when $s$ is literally $0$ and Intro selects progress,
and when Intro selects cases.  No equality congruence has been smuggled
into this step: each endpoint of the case proof is one of the two given
equations by literal normalization.

Apply Row to the resulting $D B$, opening its witness at fresh $b$.
Equations~\eqref{eq:row-three-instances} make its zero test succeed and
its first row exactly $b=t$.  Thus its selected target is $t$.  In the
single formula branch, Use of the opened clause at $2$, with reflexivity
of $0$ as antecedent, gives $F=(N(h(b))=r)$.  The exported formula is
$N(F[t/b])$, which is $N(h(t))=r$ by
\eqref{eq:row-normalization-coherence}.  The parameter $b$ is absent from
$t,r,h$ and the outer context, so the output scope condition holds.
\end{proof}

\begin{corollary}[Equality algebra and formula transport]
\label{cor:row-formula-transport}
The new calculus derives symmetry, transitivity, term congruence,
crossed one-term transport, and transport in every Eq/$D$ formula template.
\end{corollary}
\begin{proof}
For symmetry of $s=t$, use
Theorem~\ref{thm:row-term-transport} with $h(v)=v$ and $r=s$; its other
premise is reflexivity $s=s$, and its result is $t=s$.  Applying symmetry
to the conclusion of that theorem yields crossed transport
$s=t,\,h(s)=r\vdash r=h(t)$.  The usual instances now give transitivity
and term congruence.  Two crossed transports replace a parameter in
both sides of an equation.

For a compound formula template, argue externally by induction on its
finite syntax.  Suppose it is $D_{x,y}K(v,x,y)$, and its source instance
at $s$ is proved.  Open that instance at a fresh witness $c$.  To introduce
the target instance at $t$, choose witness $c$.  Since this witness is a
parameter rather than literal $0$, Intro is in case mode.  Under its
hypothesis $K(t,c,0)$, use the induction hypothesis and $t=s$ to derive
$K(s,c,0)$.  Use the source clause at $Sz$, then the induction hypothesis
and $s=t$, obtaining $K(t,c,Sz)$.  Close Intro and the source elimination.
Each recursive transport concerns the strictly smaller formula $K$.

The proof always introduces the compound target at a fresh parameter;
it does not require an unproved repair of case introduction at witness
$0$.  Its induction is external on formula syntax, not object-language
induction for quantified formulas.
\end{proof}

\subsection{Exact equivalence and its interpretation consequence}

\begin{theorem}[The entire open consequence relation is preserved]
\label{thm:row-exact-equivalence}
For every finite canonical clause context $\Gamma$ containing
$\mathcal R$ and every canonical formula $F$,
\[
 \Gamma\vdash_{\mathcal K_{\mathrm{row}}}F
 \quad\Longleftrightarrow\quad
 \Gamma\vdash_{\mathcal K_{=}}F.
\]
Both directions are effective transformations of finite proof trees and
leave the entire end sequent literally unchanged.
\end{theorem}
\begin{proof}
From the new calculus to the old, use
Corollary~\ref{cor:row-soundness} at every Row node.  The two other
instructions are unchanged.

Conversely, translate an old derivation recursively.  At an old Name
node, its translated branches prove $b=t$ and $F$ under the opened
clause.  Use Corollary~\ref{cor:row-formula-transport} inside that branch
to obtain $N(F[t/b])$.  The old rule's side condition says this result is
$b$-free, so Lemma~\ref{lem:row-ordinary-elimination} discharges the
witness.  This construction also covers targets $t$ that contain $b$;
equality transport is simultaneous one-pass substitution into the fixed
formula template.  Use and Intro are shared.

The constructions used for this replacement were derived before this
translation: ordinary elimination, then numerical cases, then the
two-row term-transport package, then formula transport.  Thus recursion
on the old proof tree is well founded.  No arbitrary numerical
substitution theorem for derivations is assumed.
\end{proof}

\begin{corollary}[Full first-order iPRA]
\label{cor:row-ipra}
The new presentation is equiconsistent and bi-interpretable with full
first-order intuitionistic primitive recursive arithmetic.  Both
numerical interpretations use the whole domain and actual equality;
both numerical composite isomorphisms have graph $x=y$.
\end{corollary}
\begin{proof}
Compose Theorem~\ref{thm:row-exact-equivalence} with the attributed
canonical/raw and native/iPRA comparisons.  The new comparison is the
identity on numerical terms and formulas and has exactly the same
ordinary reading.  It therefore transfers all proofs of the inherited
term composites, logical macros, and formula round trips without
changing their statements.  In particular, for native $F$ and ordinary
iPRA $\varphi$, the inherited statements
\[
 F\leftrightarrow_{\mathrm{native}}(F^I)^J,
 \qquad
 \varphi\leftrightarrow(\varphi^J)^I
\]
remain provable in the respective theories.  This includes the binder
round trip
\[
 D_{x,y}B\leftrightarrow_{\mathrm{native}}
 \exists_{\mathrm{native}}x\bigl(
 B(x,0)\to_{\mathrm{native}}
 \forall_{\mathrm{native}}y\,B(x,y)\bigr).
\]
No restriction to standard natural numbers, proper definable cut, or
quotient domain is introduced.  The contradiction $0=S0$ is unchanged,
so the effective translations also give equiconsistency.
\end{proof}

\paragraph{Precisely what is not claimed.}
The mode in \eqref{eq:row-target} is not invariant under arbitrary
numerical substitution: an instance can normalize to $0=0$ when its
unsubstituted zero row did not.  Consequently, this proof does not claim
that blindly substituting into a derivation preserves its primitive
instruction tags.  General admissibility follows through the exact
equivalence instead.  Nor do fewer premises imply a smaller checker
without measurement: the fixed-row normalization and tests must be
counted.  The result establishes a stricter complete equality interface,
not an absolute minimum over all possible presentations.
